\section{Conclusions and Future Works}
\label{sec5}

\noindent This paper presents MAT4PM, an active monitoring framework that integrates machine learning based control plane with a programmable data plane to enable adaptive and fine-grained threshold control in SDN. 
% The proposed system leverages the P4 language to implement lightweight monitoring logic on the switch side, in conjunction with a controller-side learning module that dynamically adjusts monitoring intervals, coordinated via a P4Runtime-based feedback interface. 
\rev{The system leverages the P4 language to implement lightweight monitoring logic on the switch side. 
A controller-side learning module dynamically adjusts monitoring intervals through a P4Runtime-based feedback interface.} 
This control interface connects the data and control planes, enabling timely synchronization between observed traffic conditions and policy updates. 
By introducing a composite cost function that jointly considers monitoring accuracy and communication overhead, MAT4PM effectively balances performance and efficiency through event-driven threshold updates.

Comprehensive evaluations in the Mininet environment and on BMv2 software switches demonstrate that the dynamic threshold schemes powered by machine learning models consistently outperform traditional fixed-threshold baselines across multiple metrics. 
Specifically, the dynamic schemes reduce average monitoring error from 9.1\% to 7.0\%, improve reporting precision, and lower overall monitoring cost by 5.6\%. 
The inference delay of the models remains below the millisecond scale, and resource consumption is minimal, indicating the practical deployability and scalability of the proposed approach in real-world network conditions.

\rev{
Future research directions include several complementary extensions to further enhance the practicality and robustness of MAT4PM. 
First, reinforcement learning and other online or self-supervised learning paradigms will be investigated to reduce or eliminate the reliance on offline labeled datasets, enabling continuous adaptation under evolving traffic conditions without manual data annotation. 
Second, deployment on physical programmable hardware platforms will be explored to assess system behavior under realistic traffic workloads and hardware constraints. 
In this context, additional aspects that remain unexplored in the current study, such as in-network security protection, fault tolerance under switch or controller failures, and energy efficiency of continuous monitoring and control operations, will be systematically examined. 
These directions aim to extend MAT4PM toward a more secure, resilient, and energy-aware intelligent monitoring framework suitable for large-scale operational networks.
}
